
\chapter{ORIUS System Context}
\label{ch:orius-context}

Grid-scale battery energy storage systems (BESS) sit at the intersection of
three economic forces: variable renewable generation that must be absorbed,
load profiles that demand dispatchable power on sub-hourly timescales, and
wholesale markets that price energy at volatile clearing rates.  This chapter
positions the ORIUS framework within that operational context, traces the
closed decision loop from forecast to measurement, and identifies the
\emph{Optimize--Dispatch boundary} as the location where observation quality
most directly threatens safety.

% ----------------------------------------------------------------
\section{Why Grid-Scale Batteries Matter}

The penetration of wind and solar generation has transformed balancing-area
operations.  In 2024 the ERCOT interconnection recorded negative real-time
prices for 18\% of daytime hours, while evening ramp rates exceeded
12\,GW\,h$^{-1}$ on 37 days.  Grid-scale batteries absorb these mismatches:
they charge when marginal cost is low, discharge when it is high, and provide
ancillary services (frequency regulation, spinning reserve) that stabilise the
grid during transients.

Three properties distinguish BESS from conventional peakers:
\begin{enumerate}
  \item \textbf{Bidirectional power flow.}  A 100\,MW / 400\,MWh installation
    can charge or discharge at full rate within seconds, making it a
    controllable load \emph{and} a controllable generator.
  \item \textbf{State-of-charge coupling.}  Every dispatch decision at time~$t$
    constrains the feasible set at time~$t+1$ through the SOC dynamics.  The
    controller cannot treat time steps independently.
  \item \textbf{Degradation risk.}  Deep cycling, high C-rates, and
    over-voltage operation accelerate cell ageing.  Safety envelopes are
    therefore both \emph{hard} (BMS trip thresholds) and \emph{soft}
    (warranty-preserving SOC margins).
\end{enumerate}

These properties mean that dispatch quality depends not only on forecast
accuracy but on the fidelity of the real-time observation that informs each
control step.

% ----------------------------------------------------------------
\section{ORIUS as a Closed Decision Loop}

The ORIUS system implements a five-stage closed loop:

\[
  \underbrace{\text{Forecast}}_{\hat{y}_{t+1:t+H}}
  \;\longrightarrow\;
  \underbrace{\text{Optimize}}_{a_t^\star}
  \;\longrightarrow\;
  \underbrace{\text{Dispatch}}_{a_t^{\mathrm{exec}}}
  \;\longrightarrow\;
  \underbrace{\text{Measure}}_{\tilde{z}_t}
  \;\longrightarrow\;
  \underbrace{\text{Monitor}}_{\mathcal{C}_t}
  \;\longrightarrow\;
  \text{(next cycle)}
\]

\begin{description}
  \item[Forecast.]  A model family $\mathcal{F}$ produces point forecasts and
    prediction intervals for load, wind, solar, and price over a rolling
    horizon $H$ (typically 24--48\,h at hourly resolution).

  \item[Optimize.]  A dispatch optimiser receives the forecast distribution
    and the current SOC estimate, then solves a constrained programme to
    produce a candidate action $a_t^\star = (\text{charge}_t,
    \text{discharge}_t)$.

  \item[Dispatch.]  The candidate action is checked against safety constraints
    (SOC bounds, ramp limits, BMS thresholds).  If the action is feasible it
    is forwarded to the battery management system; otherwise it is repaired to
    the nearest safe action $a_t^{\mathrm{safe}}$.

  \item[Measure.]  Operational telemetry $\tilde{z}_t$ returns from the plant:
    actual SOC, terminal voltage, cell temperatures, grid-side power meter
    readings.  This measurement may be delayed, dropped, corrupted, or stale.

  \item[Monitor.]  A certificate $\mathcal{C}_t$ records the action taken, the
    observation quality score, the interval width, and whether the step passed
    all safety invariants.  The certificate feeds the next forecast cycle and
    is archived for audit.
\end{description}

% ----------------------------------------------------------------
\section{The Optimize--Dispatch Boundary as the True Risk Location}

Not all stages of the loop carry equal risk.  A forecast error at
$t + 12$\,h can be corrected by re-optimisation at $t + 1$; a measurement
delay at the \emph{Measure} stage is visible to the \emph{Monitor} and can
trigger a conservative hold.  The highest-risk transition is the
\textbf{Optimize $\to$ Dispatch} boundary, for three reasons:

\begin{enumerate}
  \item \textbf{Commitment.}  Once the BMS receives a power setpoint, the
    battery executes it within milliseconds.  The action is irreversible on the
    control-cycle timescale.
  \item \textbf{Observation dependence.}  The optimiser's candidate~$a_t^\star$
    is conditioned on the most recent observation~$\tilde{z}_{t-1}$.  If that
    observation is stale or corrupted, the candidate may lie outside the true
    feasible set without any visible warning.
  \item \textbf{Constraint coupling.}  An unsafe action at time~$t$ shifts the
    SOC into a region that may render \emph{all} future actions infeasible
    until a recovery period elapses (the ``SOC trap'').
\end{enumerate}

The DC3S adapter (Chapter~\ref{ch:dc3s-adapter}) is positioned precisely at
this boundary: it intercepts~$a_t^\star$, inflates the uncertainty set to
account for observation degradation, and projects the action into the
tightened safe set before it reaches the BMS.

% ----------------------------------------------------------------
\section{Operational Telemetry Pathways}

Telemetry in a grid-scale BESS traverses multiple communication layers before
reaching the dispatch controller:

\begin{enumerate}
  \item \textbf{Cell-level sensors} (voltage, temperature, current) report to
    the Battery Management System (BMS) over CAN bus at 1--10\,Hz.
  \item \textbf{BMS aggregation} computes pack-level SOC, state-of-health
    (SOH), and thermal limits, then publishes to the site SCADA at 1\,Hz.
  \item \textbf{SCADA gateway} relays readings to the Energy Management System
    (EMS) via Modbus/TCP or IEC~61850.  This link is typically on-site LAN
    with sub-second latency.
  \item \textbf{Cloud telemetry} streams from the EMS to the dispatch
    controller (ORIUS) over cellular, fibre, or satellite uplink.  This
    link is the primary source of packet loss, delay jitter, and reordering.
  \item \textbf{Grid-side meters} (revenue meters, PMUs) provide independent
    power and frequency measurements at 4\,s or sub-second resolution.
\end{enumerate}

Each hop introduces a distinct failure mode: CAN bus saturation at high
C-rates, SCADA gateway reboot during firmware updates, cellular dropout during
storms, and satellite latency during congestion.  The telemetry fault taxonomy
used in CPSBench (Chapter~\ref{ch:cpsbench}) mirrors these real failure modes.

% ----------------------------------------------------------------
\section{Why Measurement Quality Enters the Control Problem}

Classical dispatch formulations assume that the controller observes the true
state~$z_t$ at each step.  In practice the controller observes a degraded
version~$\tilde{z}_t$ that may differ from~$z_t$ in timing, magnitude, or
completeness.  The gap between $z_t$ and $\tilde{z}_t$ has three consequences
for safety:

\begin{enumerate}
  \item \textbf{SOC estimation error.}  If the reported SOC is stale by
    $\Delta t$ steps, the true SOC has drifted by up to
    $P_{\max} \cdot \Delta t / E_{\mathrm{cap}}$ in normalised units, where
    $P_{\max}$ is the rated power and $E_{\mathrm{cap}}$ the nameplate energy.
  \item \textbf{Forecast conditioning error.}  The most recent load and
    renewables readings condition the next forecast window.  Corrupted readings
    propagate into biased intervals, undermining conformal coverage guarantees.
  \item \textbf{Constraint tightening mismatch.}  A robust optimiser that
    tightens constraints based on \emph{assumed} interval width may
    under-tighten if the true observation quality is worse than assumed, or
    over-tighten if quality is better, sacrificing economic value.
\end{enumerate}

The central thesis of this dissertation is that measurement quality must be
\emph{explicitly scored and fed back} into the uncertainty set at each control
step.  The observation quality estimator (OQE, Chapter~\ref{ch:dc3s-adapter})
and the reliability-adaptive conformal certificate (RAC-Cert) implement this
feedback loop.

% ----------------------------------------------------------------
\section{Chapter Summary}

This chapter established the ORIUS system as a five-stage closed decision
loop---Forecast, Optimize, Dispatch, Measure, Monitor---operating over
grid-scale battery storage.  The Optimize--Dispatch boundary was identified as
the highest-risk transition because actions are irreversible, observation-dependent,
and coupled through SOC dynamics.  Telemetry traverses multiple communication
layers, each introducing distinct degradation modes.  Because measurement
quality directly affects SOC estimation, forecast conditioning, and constraint
tightening, it must enter the control formulation as a first-class variable.
The remainder of Part~II develops the mathematical machinery to accomplish this.
